% Generated draft: requires mathematical and source review.
\documentclass[11pt]{article}
\usepackage[margin=25mm]{geometry}
\usepackage{fontspec}
\setmainfont{Latin Modern Roman}
\usepackage{amsmath,amssymb,mathtools,unicode-math}
\setmathfont{Latin Modern Math}
\usepackage{enumitem,needspace,xurl,hyperref}
\setlength{\parindent}{0pt}
\setlength{\parskip}{5pt}
\setlength{\emergencystretch}{3em}
\newcommand{\N}{\mathbb{N}}
\newcommand{\Z}{\mathbb{Z}}
\newcommand{\Q}{\mathbb{Q}}
\newcommand{\R}{\mathbb{R}}
\newcommand{\C}{\mathbb{C}}
\newcommand{\F}{\mathbb{F}}
\newcommand{\A}{\mathbb{A}}
\newcommand{\PP}{\mathbb{P}}
\newcommand{\E}{\mathbb{E}}
\newcommand{\eps}{\varepsilon}
\newcommand{\dd}{\,\mathrm{d}}
\begin{document}
\section*{UnsolvedMath ID: 6000005}
\textbf{Source problem code:} catalogue entry
\subsection{Definitions and mathematical statement}
In a common information-geometric convention, a statistical manifold is a smooth manifold $M$ with a Riemannian metric $g$ and a torsion-free connection $\nabla$ for which $\nabla g$ is totally symmetric. Its dual connection $\nabla^*$ satisfies
\[
 Xg(Y,Z)=g(\nabla_XY,Z)+g(Y,\nabla_X^*Z).
\]
It is dually flat when both $\nabla$ and $\nabla^*$ are flat. The source does not explicitly state whether these conventions, or a more general metric/connection class, are intended.

A regular probability-family realization would use a sample measure space $(\Omega,\lambda)$ and smooth positive densities $p_\theta$, $\theta\in M$, with $\int_\Omega p_\theta\,\mathrm d\lambda=1$, satisfying the regularity needed for the following integrals. In local coordinates put $\ell_\theta=\log p_\theta$. The Fisher metric and one standard induced connection, the exponential connection, are
\[
 g^{p}_{ij}=\int_\Omega
  (\partial_i\ell_\theta)(\partial_j\ell_\theta)p_\theta\,\mathrm d\lambda,
 \qquad
 \Gamma^{(e),p}_{ij,k}=\int_\Omega
  (\partial_i\partial_j\ell_\theta)(\partial_k\ell_\theta)
       p_\theta\,\mathrm d\lambda.
\]
Other statistical-connection conventions require separate specification.

The question is to determine necessary and sufficient conditions for a given statistical manifold to admit a family whose induced metric and intended connection coincide with its data, up to a statistical isomorphism. In the displayed convention this means $g^p=g$ and $\nabla^{(e),p}=\nabla$ after identifying parameter manifolds. The source separately asks for the analogous characterization when the given manifold is dually flat.

The sample-space class, allowed densities, global versus local realization, identifiability, induced connection, and intended meaning of ``coincides'' are not supplied. These choices are essential to the realization problem; the displayed probability model is an explicit conventional framework, not an imposed completion of the source.
\subsection{Short English statement}
Which abstract statistical manifolds arise from actual families of probability distributions with the same information-geometric structures? Determine the answer also for dually flat manifolds.

The supplied summary describes the dually flat case as essentially answered locally by exponential-family constructions and KL/canonical divergence, citing Amari--Nagaoka. It calls the general case open and says precise obstructions are not established. Its phrase ``curved or full'' exponential families and its global scope are not accompanied by a theorem or exact hypotheses, so no blanket realization statement is adopted. The registry's partial-progress label and the 1998 source's open formulation are retained with current status needing review.
\subsection{Sources}
\begin{itemize}
\item Local UnsolvedMath registry, permanent ID 6000005, source code catalogue entry, ``Statistical Manifolds from Probability Families,'' proposed by Shun-ichi Amari, 1998. Dataset curation and metadata: Ulam AI and UnsolvedMath Contributors, CC BY 4.0.
\item Furuhata, Matsuzoe, and Urakawa, \emph{Open Problems in Affine Differential Geometry and Related Topics} (1998), item 1(e), scanned PDF pages 1--2, \url{https://www.jstage.jst.go.jp/article/iis/4/2/4_2_125/_pdf/-char/en}; manual OCR from page images, accessed 2026-07-29.
\item The supplied summary cites Amari--Nagaoka for local dually flat/exponential-family constructions without a title, exact theorem, or complete bibliographic details.
\item Source-specific provenance: the article presents this as an open problem and current status NEEDS\_REVIEW. The J-STAGE article page marks the article CC BY 4.0; this source attribution is retained in addition to the dataset attribution.
\end{itemize}
\end{document}
